\begin{env}[4.1.3]
\label{0.4.1.3-fr}
No volveremos sobre la definición de los $\mathcal{A}$-\emph{Módulos} o
\emph{haces algebraicos} sobre un espacio anillado $(X,\mathcal{A})$
(G, II, 2.2). Cuando $\mathcal{A}$ es un haz de anillos no necesariamente
conmutativos, por $\mathcal{A}$-Módulo se entenderá siempre
«$\mathcal{A}$-Módulo a la izquierda», salvo indicación expresa en contrario.
Los sub-$\mathcal{A}$-Módulos de $\mathcal{A}$ se llamarán
\emph{haces de ideales} (a la izquierda, a la derecha o biláteros) en
$\mathcal{A}$ o $\mathcal{A}$-\emph{Ideales}.

Cuando $\mathcal{A}$ es un haz de anillos conmutativos, si en la definición de
los $\mathcal{A}$-Módulos se sustituye en todas partes la estructura de
\emph{módulo} por la de \emph{álgebra}, se obtiene la definición de una
$\mathcal{A}$-\emph{Álgebra} sobre $X$. Equivalentemente, una
$\mathcal{A}$-Álgebra (no necesariamente conmutativa) es un
$\mathcal{A}$-Módulo $\mathcal{C}$ provisto de un homomorfismo de
$\mathcal{A}$-Módulos
$\varphi:\mathcal{C}\otimes_{\mathcal{A}}\mathcal{C}\to\mathcal{C}$ y de una
sección $e$ sobre $X$ que satisfacen las dos condiciones siguientes.
1\textsuperscript{o} El diagrama
\oldpage[0\textsubscript{I}]{37}
\[
\begin{tikzcd}[column sep=large,row sep=large]
\mathcal{C}\otimes_{\mathcal{A}}\mathcal{C}\otimes_{\mathcal{A}}\mathcal{C}
  \arrow[r,"\varphi\otimes 1"]
  \arrow[d,"1\otimes\varphi"']
&
\mathcal{C}\otimes_{\mathcal{A}}\mathcal{C}
  \arrow[d,"\varphi"]
\\
\mathcal{C}\otimes_{\mathcal{A}}\mathcal{C}
  \arrow[r,"\varphi"']
&
\mathcal{C}
\end{tikzcd}
\]
es conmutativo. 2\textsuperscript{o} Para todo abierto $U\subset X$ y toda
sección $s\in\Gamma(U,\mathcal{C})$, se tiene
$\varphi((e|U)\otimes s)=\varphi(s\otimes(e|U))=s$. Decir que
$\mathcal{C}$ es una $\mathcal{A}$-Álgebra conmutativa equivale a decir,
además, que el diagrama
\[
\begin{tikzcd}[column sep=large,row sep=large]
\mathcal{C}\otimes_{\mathcal{A}}\mathcal{C}
  \arrow[rr,"\sigma"]
  \arrow[dr,"\varphi"']
&&
\mathcal{C}\otimes_{\mathcal{A}}\mathcal{C}
  \arrow[dl,"\varphi"]
\\
& \mathcal{C} &
\end{tikzcd}
\]
es conmutativo, donde $\sigma$ designa la simetría canónica del producto
tensorial $\mathcal{C}\otimes_{\mathcal{A}}\mathcal{C}$.

Los homomorfismos de $\mathcal{A}$-Álgebras se definen del mismo modo que los
homomorfismos de $\mathcal{A}$-Módulos en (G, II, 2.2), aunque, naturalmente,
ya no forman un grupo abeliano.

Si $\mathcal{M}$ es un sub-$\mathcal{A}$-Módulo de una
$\mathcal{A}$-Álgebra $\mathcal{C}$, la \emph{sub-$\mathcal{A}$-Álgebra de
$\mathcal{C}$ engendrada por $\mathcal{M}$} es la suma de las imágenes de los
homomorfismos $\bigotimes^n\mathcal{M}\to\mathcal{C}$ (para los
$n\geqslant 0$). Es también el haz asociado al prehaz
$U\to\mathcal{B}(U)$ de álgebras, donde $\mathcal{B}(U)$ es la subálgebra de
$\Gamma(U,\mathcal{C})$ engendrada por el submódulo
$\Gamma(U,\mathcal{M})$.
\end{env}
